%==============================================================================
%  FitQuest -- Final Internship Report
%  Milton Estuardo Torres Meraz
%  Mitacs Globalink Research Internship 2026
%  Persuasive Computing Lab, Dalhousie University
%
%  Compiles standalone with pdfLaTeX. No external images or bibliography files
%  are required.
%
%  Sections 4-9 are written as paper-ready material. Sections 2, 10 and 11 are
%  specific to the internship and would not appear in a paper.
%==============================================================================

\documentclass[11pt,a4paper]{article}

\usepackage[utf8]{inputenc}
\usepackage[T1]{fontenc}
\usepackage[english]{babel}
\usepackage[margin=2.5cm]{geometry}
\usepackage{booktabs}
\usepackage{tabularx}
\usepackage{longtable}
\usepackage{array}
\usepackage{enumitem}
\usepackage{xcolor}
\usepackage{listings}
\usepackage{titlesec}
\usepackage{fancyhdr}
\usepackage{microtype}
\usepackage[hidelinks]{hyperref}

% ---- Layout -----------------------------------------------------------------
\pagestyle{fancy}
\fancyhf{}
\fancyhead[L]{\small FitQuest --- Final Internship Report}
\fancyhead[R]{\small Torres Meraz}
\fancyfoot[C]{\small \thepage}
\renewcommand{\headrulewidth}{0.4pt}

\titleformat{\section}{\normalfont\Large\bfseries}{\thesection}{1em}{}
\titleformat{\subsection}{\normalfont\large\bfseries}{\thesubsection}{1em}{}

\setlist[itemize]{itemsep=2pt, topsep=4pt, parsep=0pt}
\setlist[enumerate]{itemsep=2pt, topsep=4pt, parsep=0pt}

% ---- Helpers ----------------------------------------------------------------
% Do not name these \partial or \0 -- LaTeX already defines those.
\newcommand{\code}[1]{\texttt{\small #1}}
\newcolumntype{L}[1]{>{\raggedright\arraybackslash}p{#1}}
\newcommand{\done}{Delivered}
\newcommand{\halfdone}{Partial}

\lstset{
  basicstyle=\ttfamily\footnotesize,
  breaklines=true,
  frame=single,
  framesep=4pt,
  rulecolor=\color{gray!50},
  backgroundcolor=\color{gray!5},
  columns=fullflexible,
  keepspaces=true,
  showstringspaces=false,
  literate={-}{{-}}1 {_}{{\_}}1
}

%==============================================================================
\begin{document}
%==============================================================================

\begin{titlepage}
\centering
\vspace*{2cm}

{\Huge\bfseries FitQuest\par}
\vspace{0.4cm}
{\LARGE A Sensor-Fusion Engine for Physical Interaction\par}

\vspace{0.8cm}
{\large Final Internship Report\par}

\vspace{1.5cm}
\rule{\textwidth}{0.5pt}
\vspace{0.6cm}

{\large\bfseries Milton Estuardo Torres Meraz\par}
\vspace{0.3cm}
{\large Mitacs Globalink Research Intern\par}
\vspace{0.2cm}
{\large Persuasive Computing Lab\\Faculty of Computer Science, Dalhousie University\par}
\vspace{0.2cm}
{\large Halifax, Nova Scotia, Canada\par}

\vspace{0.8cm}
\rule{\textwidth}{0.5pt}

\vspace{1.2cm}
\begin{tabular}{@{}ll@{}}
\textbf{Original project title} & IronQuest 3D --- Smart Dumbbell Fitness Gaming Platform \\[2pt]
\textbf{Delivered project} & FitQuest --- Sensor-Fusion Engine for Physical Interaction \\[2pt]
\textbf{Supervisors} & Dr.\ Rita Orji \\
                     & Dr.\ Fidelia Orji \\
                     & Dr.\ Grace Ataguba \\[2pt]
\textbf{Home institution} & Universidad Autónoma de San Luis Potosí, Mexico \\[2pt]
\textbf{Internship start} & May 18, 2026 \\[2pt]
\textbf{Maximum duration} & 12 weeks \\[2pt]
\textbf{Project deadline} & August 7, 2026 \\[2pt]
\textbf{Source repository} & \small\url{https://github.com/MiltonMeraz12/FitQuest---Persuasive-Computing-Lab} \\
\end{tabular}

\vfill
{\large August 2026\par}
\end{titlepage}

%==============================================================================
\tableofcontents
\newpage
%==============================================================================

\section{Executive Summary}

This report documents a twelve-week Mitacs Globalink Research Internship
conducted at the Persuasive Computing Lab, Dalhousie University. The project
began under the title \emph{IronQuest 3D --- Smart Dumbbell Fitness Gaming
Platform} and was renarrowed on June 23, 2026 to \emph{FitQuest}, a sensor-fusion
system that converts camera, inertial, and physiological measurements into a
normalized signal specification, together with a browser application that
consumes those signals to demonstrate the complete path from physical movement
to interaction.

The primary deliverable is the middleware layer rather than the application
built upon it. The system publishes a documented, per-user calibrated payload
that downstream applications consume without direct knowledge of the underlying
vision models, inertial hardware, or wireless transports.

\subsection*{Summary of outcomes}

\begin{table}[h]
\centering
\small
\begin{tabularx}{\textwidth}{@{}L{5.2cm} L{3.2cm} X@{}}
\toprule
\textbf{Dimension} & \textbf{Result} & \textbf{Evidence} \\
\midrule
Sensing modalities & 3 & Camera pose and object detection, ESP32/BNO08x IMU, Garmin Venu 3 \\
Object detector & mAP@50 $=0.919$ & 80-epoch run, 7{,}332-image merged dataset \\
Signal specification & \code{game\_control}, schema \code{2026-08-fusion-v2} & 10 documented sections, JSONL export \\
Browser application & 10 exercises, 3 session modes, 4 difficulty levels & Live client over SSE and MJPEG \\
Physical prototype & 2 printed enclosure revisions & Blender source, STL, sliced G-code \\
Automated verification & 62 tests, all passing & 10 test modules \\
Codebase & 9{,}109 lines Python; 5{,}609 lines browser client & 97 tracked files \\
Documentation & 15 technical documents, 8 dated reports & 2{,}903 lines; internal links verified \\
\bottomrule
\end{tabularx}
\end{table}

Repository measurements in this report describe the release state of
August 4, 2026. Line counts exclude generated output, model weights, and
datasets, all of which are outside version control.

\subsection*{Delivery against the original mandate}

The original invitation described sensor-enabled dumbbells driving a 3D web game,
with repetition, tempo, and range-of-motion signals supporting feedback, goals,
analytics, and exercise adherence. Table~\ref{tab:promised} compares that mandate
with the delivered system.

\begin{table}[h]
\centering
\small
\caption{Original mandate compared with the delivered system.}
\label{tab:promised}
\begin{tabularx}{\textwidth}{@{}L{4.4cm} L{2.1cm} X@{}}
\toprule
\textbf{Original element} & \textbf{Status} & \textbf{Delivered form} \\
\midrule
Sensor-enabled dumbbells & Substituted & Instrumenting individual dumbbells applies only to the weights instrumented. A glove-mounted IMU combined with camera-based dumbbell detection generalizes to any weight and provides body context that the dumbbell alone cannot supply. \\
Repetition counting & \done & Validated against camera, inertial, and wearable evidence. \\
Range-of-motion signals & \done & Per-user calibrated \code{range\_utilization}, \code{arm\_extension}, \code{height\_signal}, \code{reach\_signal}. \\
Tempo signals & \halfdone & Frame timing and inertial sample cadence are logged and tempo is derivable from \code{signal\_log}, but it is not published as a distinct signal. \\
3D web game & \done & Ten exercises with a 3D avatar, scoped as a proof of concept. \\
Feedback and goals & \done & Live coaching text, per-repetition validation, results summary. \\
Analytics and adherence & \halfdone & Frame-level JSONL export is implemented; longitudinal analysis is future work. \\
Physiological context & \done & Extension beyond the original scope: heart rate, exertion level, intensity zone. \\
\bottomrule
\end{tabularx}
\end{table}

\newpage

%==============================================================================
\section{Introduction and Context}

\subsection{Institutional setting}

The Persuasive Computing Lab at Dalhousie University studies systems intended to
influence health and well-being behaviour. The internship was framed around
developing a system for monitoring health and well-being, with the specific
proposal being IronQuest 3D: a smart-dumbbell fitness gaming platform in which
sensor-enabled dumbbells feed a browser-based 3D game, and repetition, tempo, and
range-of-motion signals support feedback, goal setting, analytics, and exercise
adherence.

Two constraints shaped the work. The internship had a twelve-week maximum from a
May 18, 2026 start date, and the expected output was a research paper due
August 7, 2026, approximately eleven weeks after the start.

\subsection{Purpose and structure of this report}

This report records what was built, in what sequence, the reasoning behind the
principal engineering decisions, the measurements obtained, and the questions
that remain open. It addresses two audiences: the supervisory team, who require
the internship narrative and an account of progress; and the subsequent research
paper, for which the technical sections are written to be transferable with
minimal revision.

Sections~\ref{sec:architecture} through~\ref{sec:results} contain the technical
material. Sections~\ref{sec:chronology} and~\ref{sec:consolidation} are specific
to the internship.

\subsection{Terminology}

\emph{IronQuest 3D} refers to the original proposal and remains the repository
name. \emph{FitQuest} refers to the delivered system and is the name used in the
browser client, the watch application, and the presentation material. Both names
denote the same project at different stages of its scope.

%==============================================================================
\section{Change of Scope}
\label{sec:scope}

\subsection{Rationale}

By the third week of the internship it became apparent that the original scope
contained a sequential dependency chain that could not be completed and validated
within the available time:

\begin{enumerate}
\item Instrumented dumbbells require designing, building, and calibrating custom
      hardware for each weight.
\item A 3D game platform requires game design, asset production, and interaction
      tuning.
\item Adherence analytics requires longitudinal data from participants, which
      depends on the first two components being operational.
\end{enumerate}

Pursuing all three in parallel would have produced three incomplete components.
The contribution with the clearest research value is the intermediate layer: the
proposition that heterogeneous, asymmetrically mounted sensors can be combined
into a single per-user calibrated specification stable enough to drive a
real-time interactive application.

\subsection{Revised scope, June 23, 2026}

The project was formally renarrowed on June 23, 2026. The pivot report recorded
five decisions:

\begin{itemize}
\item The final game is deferred and reduced to a proof-of-concept slice.
\item The vision system is fixed at body pose and dumbbell/weight detection. No
      additional vision model enters scope.
\item The abandoned keypoint-analysis runtime path is removed in full, including
      code, command-line interface entries, tests, tools, model runs, and
      documentation.
\item \code{game\_control} is reworked from a game-input structure into a general
      sensor-fusion payload.
\item Per-user auto-calibration becomes a requirement rather than a refinement.
\end{itemize}

The calibration decision is the one with the greatest bearing on the research
value of the work. Without calibration, every normalized signal encodes the body
proportions and range of motion of the individual who developed it. With
calibration, comparable physical effort maps to comparable normalized values
across users, which is the property required for the specification to transfer.

\subsection{Non-goals}

The following were declared out of scope in writing and remained so:
complex game development, multiplayer functionality, production platform
features, rule-based repetition counting from fixed joint-angle thresholds, and
any vision model beyond the two named above.

%==============================================================================
\section{System Architecture}
\label{sec:architecture}

\subsection{Design principle}

The architecture follows one rule: clients consume \code{game\_control} and do
not access raw sensor output. The browser receives normalized scalars, booleans,
and categorical labels, not bounding boxes, quaternions, or Bluetooth
characteristics.

This separation has two consequences. A rehabilitation application, an adherence
monitor, or an alternative game can be developed against the same specification
without modifying the sensing code. It also makes the system testable: the 62
automated tests exercise the specification directly, with no camera or hardware
present.

\subsection{Frame-level data flow}

Each camera frame passes through the following stages:

\begin{lstlisting}
Camera frame
  -> YOLO26 body pose (17 COCO keypoints)
  -> temporal pose smoothing (blend and occlusion hold)
  -> dumbbell/weight detector (strided, with temporal tracking)
  -> body/object side association (wrist and forearm proximity)
  -> pose-stream state
  -> auto-calibrated motion signal analysis
  -> wearable context merge
  -> inertial telemetry merge
  -> game_control payload
  -> flat signal_log record
  -> OpenCV monitor / JSONL file / browser gateway
\end{lstlisting}

\subsection{Runtime modules}

\begin{table}[h]
\centering
\small
\caption{Runtime modules and responsibilities.}
\label{tab:modules}
\begin{tabularx}{\textwidth}{@{}L{3.6cm} r L{2.6cm} X@{}}
\toprule
\textbf{Module} & \textbf{Lines} & \textbf{Layer} & \textbf{Responsibility} \\
\midrule
\code{cli.py} & 2{,}180 & Orchestration & Commands, model loading, camera lifecycle, background bridges, single-instance lock. \\
\code{ui.py} & 1{,}682 & Presentation & OpenCV monitor: skeleton, detection boxes, telemetry panels. \\
\code{sensors.py} & 1{,}363 & Adapters & Wearable JSON and inertial serial/UDP normalization; transport selection. \\
\code{motion\_analysis.py} & 1{,}023 & Analysis & Per-session calibration and normalized camera signals. \\
\code{body\_context.py} & 987 & Analysis & Object-to-limb association and temporal object tracking. \\
\code{game\_controls.py} & 677 & Specification & Assembles the payload; defines the shared inertial intensity function. \\
\code{capture\_analysis.py} & 459 & Offline & Session quality reporting from recorded JSONL. \\
\code{web\_gateway.py} & 403 & Transport & HTTP server: SSE payloads, MJPEG preview, control endpoint. \\
\code{keypoints.py} & 319 & Primitives & COCO keypoint mapping, pose smoothing, geometry, stream state. \\
\midrule
\textbf{Total} & \textbf{9{,}109} & & \\
\bottomrule
\end{tabularx}
\end{table}

\subsection{Gateway implementation}

The browser gateway uses the Python standard library \code{http.server} rather
than a web framework. The demonstration must run on a laptop without internet
access and without an additional service. Server-Sent Events over HTTP are
sufficient for a unidirectional payload stream at camera frame rate, and MJPEG
over \code{multipart/x-mixed-replace} is sufficient for the preview image. A
framework would introduce a dependency and an additional failure mode without
providing capability the demonstration requires.

The preview stream is downscaled to 640\,px at JPEG quality 72. The browser
renders it in a card approximately 290\,px wide, so encoding at full camera
resolution represented roughly threefold oversampling, which added encode time on
the publishing thread and decode time in the browser and produced observable
latency between the user's movement and the displayed image. Detection continues
to operate on the full-resolution frame, so accuracy is unaffected.

%==============================================================================
\section{Sensing Subsystems}
\label{sec:sensing}

\subsection{Computer vision}

Two YOLO26 models run per frame. The body-pose model produces 17 COCO keypoints;
a separately trained detector produces \code{dumbbell}, \code{weight}, and
\code{other} boxes.

Three mechanisms address failure modes observed during testing:

\begin{description}[leftmargin=1.2cm,style=nextline,itemsep=3pt]
\item[Pose smoothing] The pose model estimates each frame independently, so
lighting changes, motion blur, or partial occlusion can displace a wrist or elbow
for a single frame. Visible joints are blended with their recent position, and a
joint that disappears is held briefly before being excluded from analysis.
\item[Strided object inference] The dumbbell detector runs every $N$ frames
(default 2), with cached boxes reused between runs. Object position changes more
slowly than pose, so this reduces inference cost without measurable loss.
\item[Temporal object tracking] A detected box is retained for a configurable
number of missed frames and smoothed across frames. A dumbbell passing in front
of the torso is regularly lost for two or three frames; without this mechanism the
system would report the weight as released mid-repetition.
\end{description}

Accepted boxes must further satisfy per-class confidence floors, frame-area ratio
bounds, and a wrist-or-forearm proximity requirement before contributing to the
loaded state. Section~\ref{sec:results} presents the measurements that motivate
this filtering.

\subsection{Inertial subsystem}

\begin{table}[h]
\centering
\small
\caption{Inertial subsystem configuration.}
\begin{tabularx}{\textwidth}{@{}L{4.6cm} X@{}}
\toprule
\textbf{Parameter} & \textbf{Value} \\
\midrule
Microcontroller & ESP32-S3 DevKitC-1, UART USB port \\
Inertial sensor & BNO08x, I\textsuperscript{2}C address \code{0x4B} \\
Bus pins & \code{SDA = GPIO8}, \code{SCL = GPIO9} \\
Telemetry interval & 66\,ms ($\approx$15\,Hz) \\
Transports & USB serial and Wi-Fi UDP, published concurrently \\
Telemetry and discovery ports & 4210 and 4211 \\
Mount & Right forearm armband, configurable \\
\bottomrule
\end{tabularx}
\end{table}

The firmware publishes each sample over serial and UDP from a single build. An
earlier serial-only sketch was maintained alongside the wireless variant, which
introduced the possibility of divergence between two firmware files. Since the
wireless sketch is a strict superset, the serial-only variant was removed during
final consolidation.

Two implementation details address observed instability. The board calls
\code{WiFi.disconnect(true)} before each reconnection attempt, which avoids a
documented ESP32 Wi-Fi library failure pattern following a dropped connection.
The laptop broadcasts a discovery packet so that the board can reacquire the
laptop address after a network change, which was necessary because the testing
environment alternated between a phone hotspot and the campus network. The
discovery packet carries a shared token; without one, any device on the same
network could redirect the telemetry stream by transmitting its own discovery
packet.

\subsection{Wearable subsystem}

The Garmin Venu 3 provides physiological and secondary motion context. It is not
used as a primary motion sensor. Three access paths were implemented:

\begin{enumerate}
\item \textbf{Connect IQ application.} A watch application transmits heart rate,
optional RR intervals, accelerometer, gyroscope, and location fields.
\item \textbf{Cloudflare Worker relay.} A serverless endpoint backed by a D1
database receives the watch payload and serves the most recent sample.
\item \textbf{Bluetooth Low Energy fallback.} Heart rate only, from the watch's
standard broadcast mode.
\end{enumerate}

The relay addresses a specific network condition. When the laptop is connected to
the phone hotspot, the phone acts simultaneously as hotspot gateway and as the
Garmin Connect bridge, and does not route the watch request back to a hotspot
client. Connect IQ additionally requires HTTPS. A persistent HTTPS relay resolves
both conditions and replaced an earlier approach based on ephemeral tunnels,
whose URL changed on each restart and required rebuilding and re-sideloading the
watch application.

The BLE fallback is disabled by default. Operating both sources against the same
wearable file allowed a BLE missing-device state to overwrite a valid Connect IQ
sample.

\subsection{Physical enclosure}

The inertial hardware was initially assembled on a breadboard with jumper wires,
which is adequate for desk testing but unsuitable for movement capture. A
protective enclosure was modelled in Blender and printed in two revisions. The
repository contains the Blender source, the exported STL, and sliced G-code for
the printer profile used.

The enclosure fixes the sensor orientation relative to the forearm. The rotation
at which the enclosure sits on the strap determines an absolute pitch offset;
Section~\ref{sec:browser} describes how repetition validation accommodates this.

%==============================================================================
\section{Signal Specification}
\label{sec:contract}

\subsection{Normalization}

The specification uses four value types: scalars normalized to $[0,1]$, booleans,
vectors (pitch, roll, yaw, acceleration), and categorical labels such as
\code{intensity\_zone}. The control mode identifier is
\code{sensor\_fusion\_engine} and the schema version is
\code{2026-08-fusion-v2}.

\begin{table}[h]
\centering
\small
\caption{Sections of the \code{game\_control} payload.}
\begin{tabularx}{\textwidth}{@{}L{3.4cm} X@{}}
\toprule
\textbf{Section} & \textbf{Contents} \\
\midrule
\code{body\_posture} & Pose array, per-side summaries, and torso hinge. \\
\code{dumbbells} & Accepted boxes and left/right associations. \\
\code{arm\_signals} & Per-side combination of camera pose, dumbbell association, and attached hardware. \\
\code{esp32\_glove} & Orientation, motion deltas, motion intensity and state, stability index, sample cadence. \\
\code{wearable\_watch} & Heart rate, derived exertion values, and watch motion state. \\
\code{user\_state} & Aggregate normalized state: \code{exertion\_level}, \code{intensity\_zone}. \\
\code{signal\_metrics} & Calibrated camera signals and bilateral comparison. \\
\code{calibration} & Calibration state, observed bounds, and quality assessment. \\
\code{axes} & Compact normalized values for direct interface binding. \\
\code{events} & Debounced event candidates. \\
\bottomrule
\end{tabularx}
\end{table}

\subsection{Per-user calibration}

At session start the system runs a calibration window, 7\,s by default, during
which the user moves through a comfortable range. Per-side minimum and maximum
bounds are recorded for each feature, and subsequent measurements are normalized
against those bounds. Comparable physical effort therefore maps to comparable
normalized values across users of differing height, limb length, and mobility,
without model retraining or per-user code changes.

Calibration reports a quality assessment of \code{good}, \code{limited}, or
\code{insufficient}. The timer and sample-count conditions can both be satisfied
while the user has moved very little, which produces a narrow span and a
saturation-prone signal for the remainder of the session. The browser's
calibration control monitors the reported quality rather than closing after a
fixed interval.

\subsection{Exercise disambiguation signals}

A curl and an overhead triceps extension bend the elbow through a comparable
range. Elbow angle alone does not separate them; the position of the upper arm
relative to the torso does. Two signal families address this:

\begin{itemize}
\item \code{upper\_arm\_angle}, derived from the hip--shoulder--elbow angle.
      $0^\circ$ corresponds to the arm at the side; $180^\circ$ to overhead.
\item \code{torso\_hinge}, derived from the angle between vertical and the line
      joining the hip midpoint to the shoulder midpoint. Near zero when upright,
      increasing with forward hinge.
\end{itemize}

Both are published in calibrated $[0,1]$ form and in raw degrees. The browser's
repetition gates use the raw degree values. Whether an arm is near the torso or
overhead is a fixed anatomical relationship and does not require per-user
calibration; further, an exercise that holds the upper arm still, such as a curl,
leaves the calibrated span narrow early in a session, which made the calibrated
form unstable enough to trigger falsely. The calibrated forms remain in
\code{signal\_log} for cross-user comparison.

\subsection{Asymmetric hardware configuration}

The system assumes by default that the inertial glove is worn on the right hand
and the watch on the left. Both are configurable from the command line. This
asymmetry reflects the expected deployment configuration and is handled
explicitly in the specification. It also produces a limitation: with a single
glove, a grip-rotation difference such as curl against hammer curl can be
verified only on the instrumented side.

%==============================================================================
\section{Browser Application}
\label{sec:browser}

\subsection{Scope}

The browser client demonstrates that the specification drives a functioning
interaction. It is a single self-contained page served by the local gateway, with
a vendored copy of the 3D library so that it operates without internet access.

\subsection{Exercise library}

Ten exercises are implemented across five validation modes and three arm
patterns.

\begin{table}[h]
\centering
\small
\caption{Implemented exercises, validation mode, and arm pattern.}
\begin{tabularx}{\textwidth}{@{}L{5.4cm} L{2.6cm} L{2.8cm} X@{}}
\toprule
\textbf{Exercise} & \textbf{Mode} & \textbf{Arm pattern} & \textbf{Primary signal} \\
\midrule
Curl & extension & alternating & Elbow bend \\
Hammer curl & extension & alternating & Elbow bend and grip roll \\
Front raise & height & alternating & Shoulder-height ceiling \\
Lateral raise & lateral & alternating & Reach from shoulder \\
Single-arm press & height & alternating & Overhead ceiling \\
Overhead triceps extension & triceps & alternating & Upper-arm angle and glove pitch \\
Bent-over row & row & alternating & Torso hinge and arm height \\
Bilateral front hold & hold & simultaneous & Sustained height, both sides \\
Double-arm press & height & simultaneous & Overhead, both sides \\
Overhead and front hold & combo & opposed & Opposed per-side targets \\
\bottomrule
\end{tabularx}
\end{table}

Sessions run in one of three modes --- coach-led, user-selected, or randomized ---
at one of four difficulty levels that determine both the repetition target and a
strictness parameter.

\begin{table}[h]
\centering
\small
\caption{Difficulty levels.}
\begin{tabular}{@{}lrr@{}}
\toprule
\textbf{Level} & \textbf{Repetition target} & \textbf{Strictness} \\
\midrule
Beginner & 8  & 0.45 \\
Advanced & 12 & 0.62 \\
Expert   & 16 & 0.75 \\
Fit      & 20 & 0.86 \\
\bottomrule
\end{tabular}
\end{table}

\subsection{Repetition validation}

Validation separates two categories of evidence:

\begin{description}[leftmargin=1.2cm,style=nextline,itemsep=3pt]
\item[Postural preconditions] Conditions that must hold for a repetition to be
counted, evaluated from raw degree measurements: whether the upper arm is
overhead for a press, whether the torso is hinged for a row. These reject
movements that are anatomically not the prescribed exercise.
\item[Corroborating evidence] Additional sensor input that increases confidence
but cannot prevent a repetition from counting. Glove stability, watch motion
state, forearm excursion range, and grip orientation fall into this category.
\end{description}

This asymmetry follows from behaviour observed in testing. When every sensor could
reject a repetition, a single delayed or briefly stale sensor stopped counting
without any indication to the user. Under the current design no single sensor can
block a valid repetition; a missing sensor reduces only the available
corroboration.

Forearm-based validation uses the \emph{excursion} of forearm elevation rather
than its absolute value, so that the constant pitch offset introduced by
enclosure rotation (Section~\ref{sec:sensing}) cancels.

An earlier adaptive-threshold implementation tracked each side's lifetime minimum
and maximum for the current exercise and only widened. A single unusually wide or
shallow early repetition therefore skewed the threshold for the remainder of the
set. The current implementation uses a 6-second rolling window, so thresholds
follow the recent range.

\subsection{Occlusion tolerance}

A dumbbell that briefly disappears behind the torso must not invalidate the loaded
state. The client applies a graded grace period whose duration depends on the
available corroborating evidence: a declared dumbbell weight extends it, glove
motion extends it further, and postures known to be occlusion-prone extend it
further still. This complements the temporal object tracking described in
Section~\ref{sec:sensing}.

%==============================================================================
\section{Design Decisions}
\label{sec:decisions}

This section records the principal design decisions taken during the internship,
together with the alternatives considered and the reasons for rejecting them.

\begin{longtable}{@{}L{3.9cm} L{4.6cm} X@{}}
\caption{Design decisions, alternatives considered, and rationale.}\\
\toprule
\textbf{Decision} & \textbf{Alternative rejected} & \textbf{Rationale} \\
\midrule
\endfirsthead
\toprule
\textbf{Decision} & \textbf{Alternative rejected} & \textbf{Rationale} \\
\midrule
\endhead
\bottomrule
\endfoot
Glove-mounted IMU with camera dumbbell detection &
Instrumenting the dumbbells &
Instrumented weights apply only to the weights instrumented. A glove generalizes to any weight and adds body context. \\
\addlinespace
Per-user calibration window &
Fixed anatomical thresholds &
Fixed thresholds encode the proportions of whoever implemented them. Calibration is required for the specification to transfer between users. \\
\addlinespace
Raw degrees for repetition gates &
Calibrated $[0,1]$ signals &
Anatomical relationships do not require calibration, and a curl leaves the calibrated span narrow early in a session. \\
\addlinespace
Corroborating evidence cannot block a repetition &
Every sensor may reject &
A stale sensor stopped counting with no indication to the user. \\
\addlinespace
Rolling 6-second threshold window &
Lifetime minimum and maximum &
Lifetime tracking allowed one outlier repetition to skew the threshold for the set. \\
\addlinespace
Persistent Cloudflare Worker relay &
Ephemeral HTTPS tunnel &
Tunnel URLs changed on restart, requiring a watch rebuild and re-sideload each time. \\
\addlinespace
Connect IQ default, BLE opt-in &
Both wearable sources active &
Two writers on one file allowed a BLE missing-device state to overwrite a valid sample. \\
\addlinespace
Timer started from the view lifecycle &
Timer started from application \code{onStart} &
Produced a reproducible on-device crash on Venu 3 firmware 17.05 with Connect IQ 6.0.2. \\
\addlinespace
Single firmware for serial and UDP &
Separate serial and wireless sketches &
Two sketches admit divergence. The wireless build is a strict superset. \\
\addlinespace
Standard-library HTTP gateway &
Web framework or WebSockets &
Offline demonstration; SSE and MJPEG are sufficient. A framework adds a dependency and a failure mode. \\
\addlinespace
640\,px preview at quality 72 &
Full-resolution preview &
The display card is approximately 290\,px wide. Oversampling added encode and decode time and produced observable latency. \\
\addlinespace
Strided object inference &
Per-frame object inference &
Object position changes more slowly than pose; cost is reduced without measurable loss. \\
\addlinespace
Vendored 3D library &
Library loaded from a CDN &
The demonstration must operate without internet access. \\
\addlinespace
No synthetic wearable generator &
Retaining a mock heart-rate writer &
Removes the possibility of simulated physiological values entering recorded evidence. \\
\end{longtable}

%==============================================================================
\section{Results and Validation}
\label{sec:results}

\subsection{Object detector}

The dumbbell and weight detector was trained for 80 epochs on a merged dataset.

\begin{table}[h]
\centering
\small
\caption{Dataset composition and final detector metrics.}
\begin{tabular}{@{}lr@{}}
\toprule
\textbf{Quantity} & \textbf{Value} \\
\midrule
Training images & 5{,}668 \\
Validation images & 1{,}093 \\
Test images & 571 \\
Classes & \code{dumbbell}, \code{weight}, \code{other} \\
\midrule
Precision (B) & 0.919 \\
Recall (B) & 0.869 \\
mAP@50 (B) & 0.919 \\
mAP@50--95 (B) & 0.744 \\
\bottomrule
\end{tabular}
\end{table}

\subsection{Per-class performance}

\begin{table}[h]
\centering
\small
\caption{Normalized confusion matrix: fraction correctly predicted per true class.}
\begin{tabular}{@{}lr@{}}
\toprule
\textbf{True class} & \textbf{Correctly predicted} \\
\midrule
\code{weight} & 0.96 \\
\code{other} & 0.84 \\
\code{dumbbell} & 0.83 \\
\midrule
Background predicted as \code{dumbbell} & 0.72 \\
\bottomrule
\end{tabular}
\end{table}

The detector identifies dumbbells present in the frame reliably, but proposes
dumbbells in background regions at a substantially higher rate than it misses
those present. Reporting mAP@50 alone would therefore overstate performance in
live operation.

This measurement is the reason the runtime does not act on raw detections.
Before a box contributes to the loaded state it must satisfy per-class confidence
floors, frame-area ratio bounds, and a wrist-or-forearm proximity requirement.
The filtered behaviour, rather than the aggregate metric, characterizes the
system in operation. The difference between the two is itself a result: a
detector with acceptable aggregate metrics may still require substantial
geometric filtering before it can serve as an interaction signal.

\subsection{Automated verification}

\begin{table}[h]
\centering
\small
\caption{Test suite composition. All 62 tests pass.}
\begin{tabularx}{\textwidth}{@{}L{5.6cm} r X@{}}
\toprule
\textbf{Module} & \textbf{Tests} & \textbf{Coverage} \\
\midrule
\code{test\_sensor\_fusion\_payload.py} & 20 & Payload structure, inertial bridges, wearable normalization, event debouncing \\
\code{test\_cli\_detection\_defaults.py} & 10 & Command defaults and explicit-flag precedence \\
\code{test\_signal\_motion\_analysis.py} & 8 & Calibration behaviour and normalized signals \\
\code{test\_web\_gateway.py} & 6 & Client serving, SSE, control queue, path-traversal guard \\
\code{test\_garmin\_connectiq\_http\_bridge.py} & 4 & Heart-rate sanitization \\
\code{test\_object\_temporal\_tracking.py} & 4 & Detector dropout hold and box smoothing \\
\code{test\_fitquest\_worker\_pull.py} & 3 & Relay polling and atomic writes \\
\code{test\_ui\_sensor\_labels.py} & 3 & Sensor status labelling \\
\code{test\_body\_context.py} & 2 & Dumbbell-to-limb association \\
\code{test\_capture\_analysis.py} & 2 & Offline session quality reporting \\
\midrule
\textbf{Total} & \textbf{62} & \\
\bottomrule
\end{tabularx}
\end{table}

The tests exercise the signal specification rather than the hardware, so the
suite completes in approximately six seconds without a camera, board, or watch
attached. This property allowed correctness to be verified between hardware
sessions, which were limited in number.

\subsection{Delivered artefacts}

\begin{table}[h]
\centering
\small
\caption{Delivered artefacts by language.}
\begin{tabular}{@{}lrl@{}}
\toprule
\textbf{Component} & \textbf{Lines} & \textbf{Language} \\
\midrule
Runtime package & 9{,}109 & Python \\
Browser client & 5{,}609 & HTML, CSS, JavaScript \\
Documentation & 2{,}903 & Markdown \\
Inertial firmware & 375 & Arduino C++ \\
Watch application & 317 & Monkey C \\
Relay worker & 176 & JavaScript \\
\bottomrule
\end{tabular}
\end{table}

\subsection{Validation performed}

\begin{itemize}
\item Full test suite: 62 tests passing.
\item Gateway: client and vendored library served, live payloads streaming,
      removed routes returning 404.
\item Browser client: page renders without console errors; live sensor status
      updating; 3D library loaded.
\item Watch application: compiles for the Venu 3 target with Connect IQ SDK
      9.2.0.
\item Two full sessions with camera, inertial glove, and watch operating
      concurrently.
\item Five recorded capture sessions processed through the offline analysis
      pipeline.
\end{itemize}

%==============================================================================
\section{Chronology}
\label{sec:chronology}

\subsection{Phase overview}

\begin{table}[h]
\centering
\small
\caption{Internship phases.}
\begin{tabularx}{\textwidth}{@{}L{2.6cm} L{3.4cm} X@{}}
\toprule
\textbf{Period} & \textbf{Phase} & \textbf{Principal outcome} \\
\midrule
May 18 -- Jun 22 & Foundation & Vision pipeline, detector training, initial architecture \\
Jun 23 -- Jun 28 & Change of scope & Scope renarrowed; specification redesigned; abandoned path removed \\
Jun 29 -- Jul 5 & Inertial bring-up & ESP32/BNO08x validated over USB and Wi-Fi \\
Jul 6 -- Jul 12 & Wearable integration & Wearable path implemented; transports unified \\
Jul 13 -- Jul 19 & Consolidation & Watch application, relay worker, enclosure modelled \\
Jul 20 -- Jul 26 & Application & Browser client, enclosure printed, exercise library \\
Jul 27 -- Aug 3 & Correctness & Session-derived corrections; validation redesigned \\
Aug 4 -- Aug 7 & Release & Repository consolidation, evidence preservation, reporting \\
\bottomrule
\end{tabularx}
\end{table}

\subsection{Weeks 1--5: foundation (May 18 -- June 22)}

The opening phase established the vision pipeline. The base YOLO26 pose and
detection weights were obtained, the merged dumbbell dataset was assembled from
two annotated sources, and the object detector was trained through its 80-epoch
run. The repository layout, capture tooling, and offline analysis pipeline were
established during this period.

This phase also produced the assessment that led to the change of scope. With the
vision component operational, the remaining work required by the original
proposal --- custom dumbbell hardware, a complete 3D game, and longitudinal
analytics --- exceeded the time available.

\subsection{Week 6: change of scope (June 23 -- June 28)}

Scope was formally renarrowed. The abandoned keypoint-analysis path was removed in
full, \code{game\_control} was reworked from a game-input structure into a general
sensor-fusion payload, auto-calibration and normalized arm signals were added,
heart rate was mapped to \code{exertion\_level} and \code{intensity\_zone}, and
flat \code{signal\_log} records were introduced. The documentation was rewritten
around the revised deliverable.

The pivot report recorded five risks. Reviewed at the close of the internship,
four materialized in some form: the detector required substantial geometric
filtering; the inertial mounting introduced an orientation offset requiring an
excursion-based approach; the wearable path required a relay; and camera-derived
signals required validation before use. The fifth, avoiding any implication that
the game is complete, was addressed through scope definition.

\subsection{Week 7: inertial bring-up (June 29 -- July 5)}

The ESP32-S3 and BNO08x were assembled and validated. The I\textsuperscript{2}C
scanner confirmed the sensor at address \code{0x4B}, firmware was written to
publish structured telemetry, and both the USB serial and Wi-Fi UDP paths were
confirmed against the Python bridge. Five capture sessions were recorded during
this period, covering still, tilt-sweep, and motion-burst conditions, and
processed through the analysis pipeline. Those summaries recorded an effective
sample rate of approximately 15\,Hz and median pose confidence above 0.84.

The breadboard assembly was functional but insufficiently robust for movement
capture, which established the requirement for the enclosure.

\subsection{Week 8: wearable integration (July 6 -- July 12)}

The transport layer was unified so that a single launcher command listens on USB
serial and Wi-Fi concurrently, removing the need to select a transport before each
run. The BLE heart-rate bridge was added, the wearable specification was extended
to carry heart-rate contact, resting and maximum references, and optional RR
intervals, and the offline analysis was extended to summarize wearable status,
heart-rate percentiles, exertion, and intensity zones.

\subsection{Week 9: consolidation (July 13 -- July 19)}

The Connect IQ watch application was developed, which required generating a
developer key and establishing a sideloading procedure. An initial on-device
crash was traced to a timer started from the application \code{onStart} lifecycle
hook; relocating it to the view lifecycle resolved the fault. The Cloudflare
Worker relay was written and deployed, replacing the ephemeral tunnel approach.
The enclosure was modelled in Blender.

\subsection{Week 10: application (July 20 -- July 26)}

The enclosure was printed and validated. The browser client was developed from an
early prototype into a client with a prescribed exercise sequence, a 3D avatar,
and a results summary; the exercise library was extended to ten exercises; and the
Wi-Fi telemetry and wearable bridges were made more robust. A second enclosure
revision was produced in response to fit observations from the first print. A VPRI
lab visit took place on July 20, for which a project introduction was prepared.

\subsection{Week 11: correctness (July 27 -- August 3)}

Two full sessions were conducted with all three sensing modalities operating
concurrently. These identified conditions that isolated testing had not: brief
dumbbell occlusions invalidating the loaded state, a single stale sensor stopping
repetition counting, and an adaptive threshold skewed by an early outlier
repetition. Each was diagnosed and corrected; the validation design described in
Section~\ref{sec:browser} is the result. Per-exercise validation was made
distinct rather than shared across several exercises, and a diagnostic view was
added.

The corrections in this period originated in full-system sessions rather than in
component testing, which indicates that the automated suite, while effective
against regression, does not substitute for integrated trials.

\subsection{Week 12: release (August 4 -- August 7)}

Repository consolidation and final reporting, described in
Section~\ref{sec:consolidation}.

%==============================================================================
\section{Repository Consolidation}
\label{sec:consolidation}

Reproducibility being a requirement for the research output, the final week was
allocated to consolidating the repository. The review covered tracked files,
cross-file references, and content excluded by version control.

\subsection{Removals}

Superseded and unreferenced material was removed: a training launcher duplicating
logic present in the command-line interface; an unreferenced dataset preparation
script; a synthetic heart-rate generator and its sample file; a workspace
organizer; a legacy regression fixture and its route; the serial-only firmware
variant; a deprecated watch application and its sideloading test; a Blender
autosave and a superseded G-code revision; and four presentation iterations
totalling approximately 14\,MB. Tracked files decreased from 124 to 83, and
subsequently to 96 once the evidence figures described below were brought under
version control.

\subsection{Unification}

Two commands defined approximately fifty command-line options in nearly identical
blocks; these now share a single definition and differ only through an explicit
defaults table, which also documents the differences. A compatibility class whose
only output was pose liveness was folded into the keypoint module. Unused browser
code and style rules were removed.

Two defects were identified during this work. The presentation module contained a
copy of the inertial intensity function written before an angular-velocity term
was added to the runtime version; the monitor consequently under-reported
movement, reporting 0.025 where the runtime and browser reported 0.617 for the
same input. The monitor now imports the runtime function. Separately, one
command's serial port default disabled the serial transport, making an explicitly
requested serial connection inoperative. Both defects resulted from duplicated
logic updated in only one location, and neither was detectable by the test suite,
since each copy was internally consistent.

\subsection{Preservation of evidence}

Training output is excluded from version control because a single run produces
several hundred megabytes. That exclusion also covered the confusion matrix,
precision-recall curves, and per-epoch metrics for the detector, which are the
measurements this report cites and which existed on a single machine without
backup. That subset is now committed with a document recording the run that
produced it. Five capture sessions were similarly reduced from 54\,MB of video and
frames to 5.6\,KB of analysis summaries before the raw media was discarded.

The relay worker's source was committed without its deployment configuration,
which meant the service could be read but not rebuilt from the repository. A
deployment configuration was added.

%==============================================================================
\section{Limitations}
\label{sec:limitations}

\begin{description}[leftmargin=1.2cm,style=nextline,itemsep=4pt]
\item[Single-camera depth] Body posture is estimated from one 2D camera. Depth is
approximated from apparent scale rather than measured.

\item[Detector false positives] As reported in Section~\ref{sec:results}, the
detector proposes dumbbells in background regions at a high rate. System
behaviour depends on geometric filtering in addition to the detector.

\item[Single-side inertial verification] With one glove, grip-rotation
differences on the uninstrumented arm cannot be verified by any current sensor and
continue to be counted from the camera's elbow-bend signal alone.

\item[Mounting offset] Enclosure rotation on the strap determines an absolute
pitch offset. Absolute forearm elevation is therefore an approximate value, and
validation relies on excursion instead.

\item[Wearable latency] Heart rate is a context and intensity signal rather than a
control or safety signal, and samples may arrive with latency.

\item[No participant evaluation] The system has been validated technically and
exercised in full sessions by the author. It has not been evaluated with
independent participants, and no claim is made regarding usability, engagement, or
adherence.

\item[Calibration assumes cooperation] The calibration window assumes the user
moves through a representative range. Quality reporting identifies when this does
not occur but cannot correct the resulting session.

\item[Proof-of-concept application] The browser client demonstrates the sensing
path. It is not a fitness product.
\end{description}

%==============================================================================
\section{Future Work}

\begin{enumerate}
\item \textbf{Participant evaluation.} The calibration mechanism is intended to
provide cross-user transferability, and this property is currently supported by
design rather than by measurement. An evaluation with independent participants is
the most direct next step.

\item \textbf{Local detector training data.} The detector was trained on general
annotated sources. Images from the deployment setting --- the specific weights,
camera angle, lighting, and occlusion conditions --- would address the background
false-positive rate reported in Section~\ref{sec:results}.

\item \textbf{Tempo as a published signal.} Frame and sample timing are already
recorded; publishing tempo would complete the repetition, range, and tempo triad
described in the original proposal.

\item \textbf{Longitudinal analysis.} The frame-level export format was designed
for cross-session analysis and has not yet been applied over time.

\item \textbf{Bilateral inertial sensing.} A second glove would remove the
single-side verification limitation, at the cost of the asymmetric configuration
the current specification handles explicitly.

\item \textbf{Rehabilitation applications.} The signal set --- range utilization,
bilateral symmetry, stability, exertion --- corresponds to measures used in
rehabilitation assessment, which represents the most direct extension beyond
fitness applications.
\end{enumerate}

%==============================================================================
\section{Conclusion}

The internship delivered a sensor-fusion system that converts camera, inertial,
and physiological measurements into a documented, per-user calibrated signal
specification, together with a browser application that consumes that
specification to prescribe exercises, validate repetitions against multi-sensor
evidence, and provide feedback during a session.

The change of scope in week six was the decision that made completion possible.
The original proposal contained three substantial deliverables with sequential
dependencies; pursuing all three would have produced three incomplete components.
Identifying the intermediate layer as the contribution, fixing the vision scope,
and reducing the application to a proof of concept produced a plan that could be
completed and validated within the available time.

Three observations extend beyond the immediate project. Aggregate detector metrics
can overstate performance in operation: the difference between an mAP@50 of 0.919
and a background false-positive rate of 0.72 separates a detector that scores well
from one that can serve as an interaction signal without additional filtering. In
a multi-sensor system, permitting every sensor to reject an event produces less
robust behaviour than treating most sensors as corroboration, since a single stale
input should reduce confidence rather than halt the system. Finally, the
correctness defects identified in this work were found in integrated sessions
rather than in component testing, which indicates that automated testing, while
necessary to prevent regression, does not replace full-system trials.

%==============================================================================
\appendix
\section{Repository Layout}

\begin{table}[h]
\centering
\small
\begin{tabularx}{\textwidth}{@{}L{3.4cm} X@{}}
\toprule
\textbf{Path} & \textbf{Contents} \\
\midrule
\code{ironquest/} & Runtime package: interface, pipeline, sensors, analysis, specification, monitor, gateway \\
\code{web/} & Browser client and vendored 3D library \\
\code{tools/} & Wearable bridges and offline stream simulator \\
\code{tests/} & Regression suite \\
\code{firmware/} & ESP32-S3 sketches \\
\code{monkey\_c/} & Connect IQ watch application \\
\code{hardware/} & Enclosure model and print files \\
\code{cloudflare/} & Relay worker and deployment configuration \\
\code{configs/} & Training profiles \\
\code{docs/} & Technical documentation, figures, and dated reports \\
\bottomrule
\end{tabularx}
\end{table}

\section{Principal Commands}

\begin{lstlisting}
# Normal operation: camera, sensors, gateway, and browser client
.\run_ironquest.bat

# Exercise the browser client without hardware
python -m tools.simulate_game_control_stream

# Inertial telemetry check
python -m ironquest check-esp32 --port auto --seconds 10 --list-ports

# Record and analyse a session
python -m ironquest capture-motion-data --label trial --source 0 --duration 30 --video
python -m ironquest analyze-capture

# Train the object detector
python -m ironquest train-combined-dumbbell-detector --device 0

# Verification
python -m pytest tests\ -q
\end{lstlisting}

\section{Payload Excerpt}

\begin{lstlisting}
{
  "status": "ready",
  "control_mode": "sensor_fusion_engine",
  "schema_version": "2026-08-fusion-v2",
  "axes": {
    "left_arm_extension": 0.72,
    "right_arm_extension": 0.64,
    "left_height_signal": 0.81,
    "right_height_signal": 0.44,
    "symmetry_score": 0.91,
    "stability_index": 0.84,
    "exertion_level": 0.38
  },
  "arm_signals": {
    "right": {
      "pose": { "arm_extension": 0.64, "range_utilization": 0.58 },
      "hardware": {
        "esp32_glove": {
          "mounted_side": "right",
          "transport": "udp",
          "orientation_euler_deg": { "pitch": 4.1, "roll": -2.0, "yaw": 91.5 },
          "motion_intensity": 0.42,
          "motion_state": "active",
          "sample_rate_hz": 15.15,
          "stability_index": 0.84
        }
      }
    }
  },
  "user_state": {
    "heart_rate_bpm": 96,
    "exertion_level": 0.3,
    "intensity_zone": "low"
  }
}
\end{lstlisting}

\end{document}
